\documentclass[]{article}

\usepackage{amsmath}
\usepackage{multirow}
\usepackage{natbib}
\usepackage{graphicx} 


\begin{document}

\subsection{Numerical models}

We conduct a comparative analysis between two distinct dynamical systems: a deterministic Hamiltonian point-vortex flow and a canonical stochastic Lévy walk model.

\subsubsection{Point-vortex system}
The motion of a passive tracer at position $\mathbf{r} = (x, y)$ is governed by the advection in a two-dimensional velocity field $\mathbf{u}(\mathbf{r}, t)$ generated by $N$ point vortices. The velocity field is the superposition of the fields induced by each vortex:
\begin{equation}
\dot{\mathbf{r}} = \mathbf{u}(\mathbf{r}, t) = \sum_{k=1}^{N} \frac{\Gamma_k}{2\pi} \frac{\hat{\mathbf{z}} \times (\mathbf{r} - \mathbf{r}_k(t))}{|\mathbf{r} - \mathbf{r}_k(t)|^2}
\end{equation}
where $\mathbf{r}_k(t)$ and $\Gamma_k$ are the position and circulation of the $k$-th vortex, respectively, and $\hat{\mathbf{z}}$ is the unit vector normal to the plane. The vortices themselves are advected by the same velocity field, creating a self-consistent, time-dependent chaotic flow. The system is Hamiltonian, and its degree of chaoticity is controlled by the number of vortices, $N$.

\subsubsection{Lévy walk model}
As a benchmark stochastic process, we employ a standard Lévy walk model. A particle moves with a constant speed $v_0$ in a straight line for a duration $\tau$. At the end of each flight, a new direction is chosen uniformly at random from $[0, 2\pi)$. The flight times $\tau$ are drawn from a heavy-tailed probability distribution:
\begin{equation}
P(\tau) \sim \tau^{-\beta-1}
\end{equation}
for $\tau \ge \tau_{min}$, where $\beta$ is the characteristic exponent that controls the transport dynamics. For $1 < \beta < 2$, the model produces superdiffusion with a theoretical mean squared displacement exponent of $\alpha = 3 - \beta$. For $\beta > 2$, the mean flight time is finite, leading to normal diffusion with $\alpha=1$.

\subsection{Simulation datasets}

The analysis is based on two primary datasets of simulated trajectories.

For the point-vortex system, we generated trajectories for four configurations with an increasing number of vortices: $N=5, 10, 20,$ and $40$. For each configuration, the vortex circulations $\Gamma_k$ were drawn from a standard normal distribution, and their initial positions were sampled uniformly within a square domain. We then simulated the trajectories of 5 passive tracers for each of the four vortex configurations. Each trajectory was integrated for a total duration of 24.75 s with a time step of 0.05 s. All analyses were performed on these noise-free tracer trajectories.

For the Lévy walk model, we generated a benchmark dataset comprising large ensembles of trajectories for four classes of the tail exponent: $\beta = 1.2, 1.5, 1.8,$ and $2.5$. These values were chosen to span the range from strong superdiffusion ($\alpha=1.8$) to normal diffusion ($\alpha=1.0$). All walkers were simulated with a constant speed of $v_0 = 1.0$ m/s.

\subsection{Statistical analysis}

We employed a suite of statistical tools to characterize and compare the transport properties of both systems.

\subsubsection{Mean squared displacement}
The primary measure of transport is the ensemble-averaged mean squared displacement (MSD), calculated as:
\begin{equation}
\langle \Delta r^2(\Delta t) \rangle = \langle |\mathbf{r}(t+\Delta t) - \mathbf{r}(t)|^2 \rangle
\end{equation}
where the average $\langle \cdot \rangle$ is taken over all possible start times $t$ and all trajectories within a given ensemble. The anomalous diffusion exponent, $\alpha$, was determined by fitting a power law, $\langle \Delta r^2(\Delta t) \rangle \propto (\Delta t)^\alpha$, to the MSD curves on a log-log scale. The fit was performed over a range of lag times from 10\% to 60\% of the maximum trajectory duration to exclude initial ballistic effects and late-time statistical noise.

\subsubsection{Velocity autocorrelation function}
To probe the temporal memory of the dynamics, we computed the normalized velocity autocorrelation function (VACF):
\begin{equation}
C_v(\Delta t) = \frac{\langle \mathbf{v}(t) \cdot \mathbf{v}(t+\Delta t) \rangle}{\langle |\mathbf{v}(t)|^2 \rangle}
\end{equation}
Tracer velocities $\mathbf{v}(t)$ were calculated from the position data using a central finite difference scheme. The VACF measures the persistence of velocity, with its decay rate indicating the characteristic decorrelation time of the flow as experienced by the tracer.

\subsubsection{Residence time distribution}
To mechanistically investigate the origins of anomalous transport, we analyzed the statistics of trapping events. A tracer was defined as being in a "trapped" state whenever its instantaneous speed fell below the 25th percentile of the pooled speed distribution for its respective N-vortex configuration. We measured the duration, $\tau_{res}$, of each continuous trapping event and constructed the probability distribution $P(\tau_{res})$. The tail of this distribution was fitted with a power law, $P(\tau_{res}) \sim \tau_{res}^{-\gamma}$, to identify Lévy-like trapping behavior.

\subsubsection{Displacement probability density functions}
We examined the evolution of the one-dimensional displacement probability density function (PDF), $P(\Delta x)$, for various lag times $\Delta t$, where $\Delta x(t, \Delta t) = x(t+\Delta t) - x(t)$. The shape of these distributions provides insight into the nature of the underlying random process. For the most superdiffusive cases (N=40 vortex system and $\beta=1.5$ Lévy walk), the empirical PDFs at long lag times were compared to both Gaussian and symmetric Lévy-stable distributions to test for convergence to a non-Gaussian attractor.

\subsubsection{Ergodicity analysis}
To quantify the statistical heterogeneity across trajectories, we measured the ergodicity-breaking (EB) parameter. First, we computed the time-averaged MSD (TAMSD) for each individual trajectory $i$:
\begin{equation}
\overline{\delta_i^2(\Delta t)} = \frac{1}{T-\Delta t} \int_0^{T-\Delta t} |\mathbf{r}_i(t'+\Delta t) - \mathbf{r}_i(t')|^2 dt'
\end{equation}
where $T$ is the total observation time. The EB parameter is then defined as the normalized variance of the TAMSDs across the ensemble of trajectories:
\begin{equation}
EB(\Delta t) = \frac{\text{Var}[\overline{\delta^2(\Delta t)}]}{(\text{Mean}[\overline{\delta^2(\Delta t)}])^2} = \frac{\langle (\overline{\delta^2})^2 \rangle - \langle \overline{\delta^2} \rangle^2}{\langle \overline{\delta^2} \rangle^2}
\end{equation}
An EB value of zero signifies perfect ergodicity, while a sustained positive value indicates that individual trajectories exhibit statistically distinct diffusive properties. We characterized the degree of non-ergodicity by calculating a plateau value, defined as the mean of $EB(\Delta t)$ over the final 20\% of the available lag times.

\end{document}
                